% !TeX root = ../../main.tex

\section{Introducción}

HDFS es el sistema de almacenamiento distribuido utilizado por Apache Hadoop.
Su función es guardar archivos muy grandes en un conjunto de computadoras y
mantenerlos disponibles incluso cuando alguno de los equipos presenta una
falla.

\section{¿Qué es HDFS?}

HDFS significa \textit{Hadoop Distributed File System}. Es un sistema de
archivos distribuido que divide los archivos en bloques y los almacena en
diferentes computadoras de un clúster.

HDFS está diseñado para almacenar grandes cantidades de información y ofrecer
un acceso de alto rendimiento. Su modelo se orienta a escribir los datos una
vez y leerlos muchas veces.

\section{Objetivos de su diseño}

Los principales objetivos del diseño de HDFS son:

\begin{itemize}
    \item Almacenar archivos de gran tamaño.
    \item Utilizar varias computadoras como un solo sistema de almacenamiento.
    \item Funcionar con hardware convencional.
    \item Mantener los datos disponibles cuando falla un nodo.
    \item Proporcionar un alto rendimiento en la lectura de datos.
    \item Facilitar el crecimiento del sistema mediante nuevos DataNodes.
    \item Ejecutar el procesamiento cerca de donde se encuentran los datos.
\end{itemize}

HDFS prioriza el rendimiento general y la tolerancia a fallos sobre el acceso
inmediato a registros individuales.

\section{Historia de HDFS}

El origen de HDFS está relacionado con Nutch, un motor de búsqueda de código
abierto desarrollado por Doug Cutting y Mike Cafarella. El crecimiento de la
información de Internet hizo necesario crear un sistema de almacenamiento
distribuido.

En 2003, Google publicó un artículo sobre Google File System. Sus ideas
sirvieron como base para desarrollar Nutch Distributed File System. Más tarde,
este sistema evolucionó y pasó a formar parte de Hadoop con el nombre de HDFS.

HDFS se convirtió en uno de los componentes principales de Apache Hadoop.
Versiones posteriores incorporaron mejoras como alta disponibilidad,
federación de NameNodes y codificación de borrado para aprovechar mejor el
espacio de almacenamiento.

\section{Características principales}

Las características principales de HDFS son:

\begin{itemize}
    \item \textbf{Almacenamiento distribuido:} los archivos se guardan en
    diferentes nodos.

    \item \textbf{Bloques grandes:} los archivos se dividen en bloques de gran
    tamaño, normalmente de 128 MB, aunque este valor puede configurarse.

    \item \textbf{Replicación:} cada bloque puede almacenarse en varios
    DataNodes.

    \item \textbf{Tolerancia a fallos:} si un nodo falla, los datos pueden
    obtenerse desde otra réplica.

    \item \textbf{Escalabilidad:} permite agregar DataNodes para aumentar la
    capacidad.

    \item \textbf{Alto rendimiento:} está optimizado para leer grandes
    cantidades de datos de manera continua.

    \item \textbf{Data Locality:} intenta ejecutar el procesamiento cerca de
    los datos.

    \item \textbf{Modelo de escritura simple:} está orientado a escribir una
    vez y leer muchas veces.
\end{itemize}

\section{Arquitectura de HDFS}

HDFS utiliza una arquitectura maestro-trabajador. El NameNode funciona como
el nodo principal y los DataNodes funcionan como nodos de almacenamiento.

El NameNode administra el espacio de nombres y los metadatos. Por ejemplo,
conoce los nombres de los archivos, sus permisos, los bloques que los forman
y los DataNodes donde se encuentran.

Los DataNodes almacenan físicamente los bloques y atienden las solicitudes de
lectura y escritura de los clientes. El cliente primero consulta al NameNode
para localizar los bloques y después se comunica directamente con los
DataNodes.

Esta separación permite que el NameNode coordine el sistema sin transportar
el contenido completo de cada archivo.

\section{Componentes principales}

Los principales componentes de HDFS son:

\begin{enumerate}
    \item \textbf{NameNode:} administra los metadatos, el espacio de nombres
    y la ubicación de los bloques.

    \item \textbf{DataNode:} almacena los bloques, atiende solicitudes de los
    clientes y ejecuta las órdenes de creación, eliminación y replicación.

    \item \textbf{Cliente HDFS:} permite crear, leer, escribir y eliminar
    archivos. Consulta al NameNode y transfiere los datos con los DataNodes.

    \item \textbf{NameNode activo y en espera:} en una configuración de alta
    disponibilidad, uno atiende las solicitudes y otro mantiene una copia
    actualizada para sustituirlo si ocurre una falla.

    \item \textbf{JournalNodes:} conservan los registros de cambios que
    permiten sincronizar los NameNodes activo y en espera.
\end{enumerate}

El Secondary NameNode crea puntos de control al combinar la imagen del sistema
de archivos con el registro de cambios. No debe confundirse con un reemplazo
automático del NameNode principal.

\section{Arquitectura maestro-esclavo}

Una arquitectura maestro-esclavo, también llamada actualmente
maestro-trabajador, es un modelo donde un componente principal coordina a
varios componentes secundarios.

En HDFS, el NameNode es el maestro porque administra los metadatos y decide
dónde deben almacenarse los bloques. Los DataNodes son los trabajadores porque
guardan los bloques y ejecutan las operaciones solicitadas.

Este modelo facilita la administración del clúster, pero también hace necesario
proteger al NameNode mediante mecanismos de alta disponibilidad.

\section{Almacenamiento por bloques}

HDFS no almacena necesariamente un archivo completo en una sola computadora.
Primero lo divide en bloques de tamaño fijo. Después distribuye esos bloques
entre diferentes DataNodes.

Por ejemplo, si el tamaño de bloque es de 128 MB, un archivo de 300 MB puede
dividirse en dos bloques de 128 MB y uno de 44 MB. Cada bloque puede tener
varias réplicas.

\subsection{NameNode}

El NameNode no almacena normalmente el contenido de los archivos. Conserva
información como:

\begin{itemize}
    \item Nombre y ruta de cada archivo.
    \item Permisos y propietario.
    \item Bloques que forman el archivo.
    \item Ubicación de las réplicas.
    \item Estado de los DataNodes.
\end{itemize}

\subsection{DataNode}

Los DataNodes almacenan los bloques en sus discos. También realizan las
siguientes actividades:

\begin{itemize}
    \item Leer y escribir bloques.
    \item Crear y eliminar bloques.
    \item Enviar reportes de bloques al NameNode.
    \item Enviar señales periódicas de funcionamiento.
    \item Replicar bloques cuando lo solicita el NameNode.
\end{itemize}

\section{Funcionamiento de HDFS}

\subsection{Proceso de lectura}

El proceso de lectura funciona de la siguiente manera:

\begin{enumerate}
    \item El cliente solicita abrir un archivo.
    \item El NameNode verifica la solicitud y localiza sus bloques.
    \item El NameNode devuelve una lista de DataNodes que poseen las réplicas.
    \item El cliente selecciona el DataNode más cercano.
    \item El DataNode envía el primer bloque directamente al cliente.
    \item El cliente continúa con los siguientes bloques hasta completar el
    archivo.
    \item Al finalizar, el cliente cierra la conexión.
\end{enumerate}

Si un DataNode no responde o entrega un bloque dañado, el cliente intenta leer
otra réplica.

\subsection{Proceso de escritura}

El proceso de escritura se realiza así:

\begin{enumerate}
    \item El cliente solicita crear un archivo.
    \item El NameNode comprueba los permisos y verifica que el archivo no
    exista.
    \item El NameNode selecciona los DataNodes donde se almacenarán las
    réplicas.
    \item El cliente divide la información en paquetes y los envía al primer
    DataNode.
    \item El primer DataNode almacena los paquetes y los envía al segundo.
    \item El segundo los almacena y los envía al siguiente DataNode.
    \item Los DataNodes devuelven confirmaciones en sentido contrario.
    \item El proceso se repite con los demás bloques.
    \item El cliente cierra el archivo y el NameNode registra la operación.
\end{enumerate}

Esta cadena de DataNodes se conoce como tubería o
\textit{replication pipeline}.

\section{Ventajas y desventajas}

\subsection{Ventajas}

\begin{itemize}
    \item Permite almacenar archivos muy grandes.
    \item Puede crecer agregando nuevos DataNodes.
    \item Mantiene copias de los bloques para evitar pérdidas.
    \item Puede utilizar hardware convencional.
    \item Ofrece buen rendimiento en lecturas secuenciales.
    \item Facilita el procesamiento paralelo.
    \item Detecta fallas y recupera réplicas automáticamente.
\end{itemize}

\subsection{Desventajas}

\begin{itemize}
    \item No es adecuado para una gran cantidad de archivos pequeños.
    \item No está diseñado para modificar continuamente partes de un archivo.
    \item La replicación consume espacio adicional.
    \item No es ideal para aplicaciones transaccionales.
    \item El acceso puede tener mayor latencia que un sistema local.
    \item Su instalación y administración requieren conocimientos técnicos.
    \item El NameNode necesita protección especial porque administra todos los
    metadatos.
\end{itemize}

\section{Replicación de datos}

La replicación consiste en guardar varias copias de cada bloque en diferentes
DataNodes. El número de copias se conoce como factor de replicación y puede
configurarse para cada archivo. Un valor común es tres.

Durante una escritura, el NameNode selecciona una tubería de DataNodes. El
primer DataNode recibe el bloque, lo guarda y lo envía al siguiente hasta
completar las réplicas.

HDFS también utiliza conocimiento de los racks. Procura colocar las copias en
más de un rack para que los datos continúen disponibles si falla un equipo,
un conmutador o un rack completo.

Si el NameNode detecta que un bloque tiene menos réplicas de las necesarias,
ordena crear una copia nueva desde alguna réplica disponible.

\section{Data Locality}

Data Locality significa acercar el procesamiento a la ubicación de los datos.
Mover una tarea pequeña suele ser más eficiente que transferir un archivo de
gran tamaño por la red.

HDFS y Hadoop utilizan tres niveles principales:

\begin{enumerate}
    \item \textbf{Node Local:} la tarea se ejecuta en el mismo DataNode donde
    está el bloque.

    \item \textbf{Rack Local:} la tarea se ejecuta en otro nodo del mismo rack.

    \item \textbf{Off-Rack:} la tarea se ejecuta en un nodo ubicado en otro
    rack.
\end{enumerate}

La primera opción es la más eficiente porque reduce el tráfico de red y mejora
el rendimiento.

\section{Manejo de fallas}

Los DataNodes envían periódicamente una señal llamada \textit{heartbeat} y un
reporte de sus bloques al NameNode. Si estas señales dejan de llegar, el
NameNode considera que el DataNode no está disponible.

HDFS maneja las fallas de las siguientes maneras:

\begin{itemize}
    \item \textbf{Falla durante una lectura:} el cliente abandona el DataNode
    con problemas y solicita el bloque a otra réplica.

    \item \textbf{Falla durante una escritura:} el DataNode defectuoso se
    elimina de la tubería y la escritura continúa con los nodos restantes.
    Después se crea la réplica faltante.

    \item \textbf{Falla de un DataNode:} el NameNode identifica los bloques
    que perdieron una copia y ordena su replicación en otros nodos.

    \item \textbf{Bloque dañado:} las sumas de verificación permiten detectar
    corrupción. El sistema utiliza una réplica válida y sustituye la copia
    dañada.

    \item \textbf{Falla del NameNode:} una configuración de alta disponibilidad
    utiliza un NameNode activo y otro en espera. Los JournalNodes mantienen
    sincronizados sus registros para permitir la sustitución.

    \item \textbf{Falla de un rack:} las réplicas colocadas en otros racks
    mantienen los datos disponibles.
\end{itemize}

Gracias a estos mecanismos, HDFS puede continuar trabajando ante fallas
comunes sin perder inmediatamente la información.

\section{Conclusión}

HDFS es un sistema de archivos distribuido diseñado para almacenar grandes
volúmenes de datos. Su arquitectura utiliza un NameNode para administrar los
metadatos y varios DataNodes para guardar bloques.

La replicación, el manejo de fallas y la ubicación local de las tareas permiten
que el sistema sea escalable y confiable. Sin embargo, HDFS funciona mejor con
archivos grandes y lecturas secuenciales que con archivos pequeños o
aplicaciones transaccionales.